\scanpage{36}
\mspage{39 bis}
\margin{\figure{\begin{tikzpicture}[x=1cm,y=1cm]\draw[gray,line width=.2pt] (0,1) -- (1.2,1.45);\draw[gray,line width=.2pt] (0,0) .. controls (.7,-.08) and (1.2,-.16) .. (1.55,-.09);\end{tikzpicture}}}
fini ($\varepsilon$ et $m_o$ fixés). Il en résulte aussitôt que la suite $\psi(m)$ est de Cauchy, donc convergente.

Soit $a$ la limite de cette suite. En remplaçant la suite donnée $\varphi(m)$ par la suite $\varphi(m)-am$, on se ramène au cas où cette limite $a$ est nulle; autrement dit, il nous suffit de prouver que, dans ce dernier cas, on a nécessairement $|\varphi(m)|\leq d$ pour tout $m$. Or on a, d'après la relation (14), $|\varphi(2m)-2\varphi(m)|\leq d$; plus généralement, pour tout entier $h>0$, on en déduit, par récurrence sur $h$, $|\varphi(2^hm)-2^h\varphi(m)|\leq(2^h-1)d$. Donc, s'il existait un $m_1$ tel que $\varphi(m_1)=d_1>d$, on aurait $\varphi(2^hm_1)>(2^h-1)(d_1-d)$, et $\psi(2^hm_1)=\varphi(2^hm_1)/2^hm_1$ ne tendrait pas vers $0$ pour $h$ infini. Ceci contredirait l'hypothèse $a=0$.
\mspage{35}
\figure{\begin{tikzpicture}\draw (0,0)--(0.02,-0.22);\end{tikzpicture}}
